\subsection{Encryption ও Decryption-এর দ্রুত উদাহরণ}

\par\smallskip\noindent
\begin{minipage}[t]{.47\linewidth}\vspace{0pt}\small
\begin{tabularx}{\linewidth}{|Y|Y|Y|}\hline
\chSixHeader Plaintext & Rule & Ciphertext\\\hline
HELLO & ৩ ঘর সামনে & KHOOR\\\hline
HELLO & ১ ঘর সামনে & IFMMP\\\hline
XYZ & ৩ ঘর সামনে & ABC\\\hline
HSC & ২ ঘর সামনে & JUE\\\hline
\end{tabularx}
\end{minipage}\hfill
\begin{minipage}[t]{.49\linewidth}\vspace{0pt}\small
\textbf{উল্টো পথে ডিক্রিপশন।} তিন ঘর Shift-এ HELLO থেকে KHOOR হয়।
এবার প্রতিটি অক্ষর তিন ঘর পিছিয়ে পাই—
\[
K\rightarrow H,\quad H\rightarrow E,\quad O\rightarrow L,
\]
\[
O\rightarrow L,\quad R\rightarrow O.
\]
ফলে আবার মূল HELLO পাওয়া যায়।
\end{minipage}\par\smallskip

\begin{keypointbox}{সরণের নিয়ম}
প্রশ্নে দেওয়া Shift বা Rule-ই ব্যবহার করবে; আগের Ciphertext মুখস্থ করে
লিখবে না। সারণিতে একই HELLO-র জন্য $k=3$ ও $k=1$-এ ভিন্ন ফল হয়েছে।
\end{keypointbox}


% ============================================================
% 16.11 CUSTOM RULE
% ============================================================

